\chapter{Probabilistic model: supporting information}
\label{app:pm}


\section{Probabilistic model}\label{app:prob}

This appendix supports Section~\ref{sec:prob}. It gives the equations of the
two limit states, the estimation of the parameters, and the checks of the
model against data and against changes in the dependence assumptions. The
model is \texttt{scripts/StressCaclulation/FailureProbability.py}.

\subsection{Load on the wall}\label{app:prob-load}

All moments are taken about the front pile row at the underside of the floor,
$z_b$ ($-0.55$\,m NAP for type 1, $-0.22$\,m NAP for type 2, plus $\Delta z$).
The loaded height is $H = z_\text{werf} - z_b$. The groundwater and the canal
stand $h_w$ and $h_c$ above $z_b$ (between 0 and $H$). The load moment is
\begin{equation}
S = \underbrace{\tfrac{1}{6}K\left(\gamma H^3 - \gamma_w h_w^3\right)}_{\text{soil, effective weight}}
  + \underbrace{\tfrac{1}{6}\gamma_w\left(h_w^3 - h_c^3\right)}_{\text{groundwater} - \text{canal}}
  + \underbrace{\tfrac{1}{2}K q_\text{eff} H^2}_{\text{surcharge}},
\qquad q_\text{eff} = q\,\frac{H^2}{H^2 + a^2},
\label{eq:app-load}
\end{equation}
and the horizontal force is
$F_h = \tfrac12 K(\gamma H^2 - \gamma_w h_w^2) + K q_\text{eff} H
+ \tfrac12\gamma_w(h_w^2 - h_c^2)$. The same $K$ and the same groundwater
level enter every term.

\subsection{Resistance against overturning}\label{app:prob-resistance}

The body that tips is the wall, its floor and the soil on the floor. Its
weight $V$ and moment $M$ about the front pile row follow from the areas of the
parts in the measured cross-section (masonry, tamped concrete, reinforced
concrete or timber, and soil). Below the mean of the canal and the groundwater
level each part weighs $\gamma - \gamma_w$. The soil behind the body drags it
down by wall friction:
\begin{equation}
M_R = \theta_R M + \tan\delta\,F_s\,(x_\text{back} - x_\text{front}),
\qquad \delta = \tfrac23\varphi, \quad \sin\varphi = \frac{1-K}{1+K},
\end{equation}
with $F_s$ the horizontal force of soil and surcharge. The wall overturns when
$M_R < S$.

\subsection{Forces in the piles}\label{app:prob-piles}

The rear row carries $R = (M_R - S)/(x_\text{rear} - x_\text{front})$; the
front row carries the rest of the weight, $F = V + \tan\delta\,F_s - R$. For
type 2 the front row is raked 1:10 and takes $0.1F$ horizontally. The rest of
the horizontal force is shared by both rows, so per pile (spacing $s$)
\begin{equation}
P = \tfrac12 (F_h - 0.1F)\,s, \qquad N = F s \sqrt{1 + 0.1^2}.
\end{equation}
The moment in a long pile on elastic springs (Hetényi) is
\begin{equation}
M = \frac{P}{2\lambda} \;\text{(fixed head, type 2)}, \qquad
M = \frac{0.3224\,P}{\lambda} \;\text{(pinned head, type 1)},
\end{equation}
with $\lambda = \left(k D_0 / 4 E I\right)^{1/4}$ and $I = \pi d^4/64$.
A sound core thinner than 20\,mm counts as broken.
The pile breaks when
\begin{equation}
\theta_M \frac{32 M}{\pi d^3} + \frac{4N}{\pi d^2} > \text{MOR},
\qquad d = D_0 - 2\bigl(s_0 + r\,(t - 2026)\bigr),
\end{equation}
where $s_0$ is today's soft shell and $r$ its growth rate.

\subsection{Parameter estimation}\label{app:prob-estimation}

\paragraph{Moments to distribution.} A lognormal with mean $m$ and CoV $v$ has
$\sigma_{\ln}^2 = \ln(1 + v^2)$ and $\mu_{\ln} = \ln m - \sigma_{\ln}^2/2$.
Normal and uniform distributions are set by their mean and CoV, or by their
bounds.

\paragraph{Proven strength.} Let $E$ be the event that the wall stands today
under soil and water, without surcharge. Every draw $\mathbf{x}_i$ is checked
against $E$ with its own resistance and soil and a separate groundwater draw.
Bayes' rule gives $f(\mathbf{x} \mid E) \propto f(\mathbf{x})\,
\mathbf{1}[\mathbf{x} \in E]$, so the updated failure probability is
estimated by rejection:
\begin{equation}
\hat P_{f\mid E} = \frac{\sum_i \mathbf{1}[\text{fail}_i]\,\mathbf{1}[E_i]}{\sum_i \mathbf{1}[E_i]}.
\end{equation}

\paragraph{Wall type.} With prior $p = P(T = 2)$ per rak, the same evidence
updates the type, and the best estimate weighs the two types:
\begin{equation}
P(T = 2 \mid E) = \frac{p\,P(E \mid T=2)}{p\,P(E \mid T=2) + (1-p)\,P(E \mid T=1)},
\qquad
P_f = \sum_{T} P(T \mid E)\,P_f(T \mid E).
\end{equation}

\subsection{Results per wall}\label{app:prob-results}

Table~\ref{tab:prob-results} gives the result of the run of 30 October 2026
(400{,}000 draws per wall and type) for every modelled wall. The prior
$P(T{=}2)$ is the chance of type~2 for the rak; the updated value follows from
Equation~(\ref{eq:mixture}). All $\beta$ values are updated with proven
strength. Kromme Nieuwe Gracht 005 to 007 get the same result, because nothing
in the data tells them apart.

\begin{table}[htbp]
\centering
\small
\caption{Reliability index $\beta$ per modelled wall, updated with proven
strength. OG = Oude Gracht, KNG = Kromme Nieuwe Gracht. Type 1 and type 2:
$\beta$ in 2026 if the wall is of that type. Best estimate: both types
weighted by $P(T \mid E)$.}
\label{tab:prob-results}
\setlength{\tabcolsep}{5pt}
\begin{tabular}{@{}l r l c r r r r@{}}
\toprule
 & & & $P(T{=}2)$ & \multicolumn{2}{c}{$\beta$ 2026, per type} & \multicolumn{2}{c}{$\beta$ best estimate} \\
\cmidrule(lr){5-6}\cmidrule(l){7-8}
Wall & Length [m] & Soil class & prior $\rightarrow$ updated & Type 1 & Type 2 & 2026 & 2041 \\
\midrule
OG Rak 03 OZ    & 101.1 & clay/peat & 0.2 $\rightarrow$ 0.34 & $-0.70$ & 1.62 & $-0.06$ & $-0.08$ \\
OG Rak 04 OZ    &  67.4 & sand      & 0.2 $\rightarrow$ 0.21 & 0.64    & 2.04 & 0.81    & 0.76 \\
OG Rak 05 NZ    &  72.1 & sand      & 0.2 $\rightarrow$ 0.21 & 0.66    & 2.05 & 0.82    & 0.78 \\
OG Rak 06 NZ    &  43.7 & sand      & 0.2 $\rightarrow$ 0.20 & 1.05    & 2.45 & 1.18    & 1.14 \\
OG Rak 06 ZZ    &  57.1 & sand      & 0.2 $\rightarrow$ 0.21 & 0.66    & 2.06 & 0.83    & 0.78 \\
OG Rak 08 OZ    &  32.7 & mixed     & 0.8 $\rightarrow$ 0.82 & 0.27    & 2.11 & 1.38    & 1.31 \\
OG Rak 09 ZOZ   &  54.5 & clay/peat & 0.8 $\rightarrow$ 0.90 & $-0.79$ & 1.50 & 1.10    & 1.00 \\
OG Rak 09 WZ    &  53.4 & clay/peat & 0.8 $\rightarrow$ 0.92 & $-0.91$ & 1.35 & 1.05    & 0.93 \\
OG Rak 09 WZ 01 &   3.4 & clay/peat & 0.8 $\rightarrow$ 0.89 & $-0.69$ & 1.65 & 1.13    & 1.04 \\
KNG 002         &  16.2 & clay/peat & 0.9 $\rightarrow$ 0.94 & $-0.65$ & 1.68 & 1.37    & 1.26 \\
KNG 003         &   9.9 & mixed     & 0.9 $\rightarrow$ 0.92 & 0.10    & 1.89 & 1.52    & 1.40 \\
KNG 004         &  11.1 & mixed     & 0.9 $\rightarrow$ 0.93 & $-0.09$ & 1.66 & 1.39    & 1.24 \\
KNG 005         &  13.3 & clay/peat & 0.5 $\rightarrow$ 0.74 & $-0.79$ & 1.27 & 0.59    & 0.52 \\
KNG 006         &   4.1 & clay/peat & 0.5 $\rightarrow$ 0.74 & $-0.79$ & 1.27 & 0.59    & 0.52 \\
KNG 007         &  10.3 & clay/peat & 0.5 $\rightarrow$ 0.74 & $-0.79$ & 1.27 & 0.59    & 0.52 \\
KNG 008         &  42.1 & clay/peat & 0.5 $\rightarrow$ 0.60 & $-0.47$ & 1.95 & 0.56    & 0.53 \\
KNG 009         &  18.3 & sand      & 0.5 $\rightarrow$ 0.51 & 1.07    & 2.46 & 1.45    & 1.40 \\
\bottomrule
\end{tabular}
\end{table}

\subsection{Sensitivity to dependence and inputs}\label{app:prob-sensitivity}

The script \texttt{SensitivityDependence.py} (in
\texttt{scripts/StressCaclulation/}) reruns the model with one change at a
time, with 200{,}000 draws per wall, type and variant and the same seeds. It
writes \texttt{walmuur\_sensitivity.csv} in \texttt{data/processed/}. The
correlations use a Gaussian copula: the two variables are drawn from
correlated standard normals and transformed to their own distributions. For
the uncertain soil class, every class found within 150\,m gets a prior weight
equal to its share of the borings; like the wall type, it is then updated
with proven strength. Table~\ref{tab:sens-summary} summarises the change in
$\beta$ (updated, 2026); Table~\ref{tab:sens-walls} gives every wall.

\begin{table}[htbp]
\centering
\small
\caption{Change in $\beta$ against the base model: median [range] over the
17 walls. Best estimate: the two wall types weighted by $P(T \mid E)$.}
\label{tab:sens-summary}
\setlength{\tabcolsep}{4pt}
\begin{tabular}{@{}l r r r@{}}
\toprule
Change & Best estimate & Type 1 & Type 2 \\
\midrule
$\rho(K,\gamma) = -0.5$ & $0.00$ [$-0.01$, $0.00$] & $+0.01$ [$-0.01$, $+0.01$] & $0.00$ [$-0.01$, $+0.01$] \\
$\rho(\text{MOR},E) = +0.8$ & $0.00$ [$0.00$, $+0.01$] & $0.00$ [$-0.01$, $+0.01$] & $+0.01$ [$0.00$, $+0.03$] \\
Groundwater +0.6 to +1.0\,m NAP & $+0.02$ [$-0.38$, $+0.08$] & $-0.22$ [$-0.43$, $-0.18$] & $-0.06$ [$-0.10$, $-0.05$] \\
Soil class uncertain & $+0.06$ [$-0.70$, $+0.56$] & -- & -- \\
\bottomrule
\end{tabular}
\end{table}

\begin{table}[htbp]
\centering
\small
\caption{$\beta$ per wall (updated, 2026, 200{,}000 draws) for each change.
Shares c/s/m: clay/peat, sand and mixed among the borings within 150\,m.
OG = Oude Gracht, KNG = Kromme Nieuwe Gracht.}
\label{tab:sens-walls}
\setlength{\tabcolsep}{4pt}
\begin{tabular}{@{}l l c r r r r r@{}}
\toprule
Wall & Class & Shares c/s/m & Base & $\rho(K,\gamma)$ & $\rho(\text{MOR},E)$ & Groundw. & Soil unc. \\
\midrule
OG Rak 03 OZ & clay/peat & 0.50 / 0.50 / 0.00 & $-0.05$ & $-0.05$ & $-0.05$ & 0.00 & 0.51 \\
OG Rak 04 OZ & sand & 0.50 / 0.33 / 0.17 & 0.81 & 0.81 & 0.81 & 0.55 & 0.38 \\
OG Rak 05 NZ & sand & 0.40 / 0.60 / 0.00 & 0.83 & 0.82 & 0.83 & 0.56 & 0.56 \\
OG Rak 06 NZ & sand & 0.50 / 0.33 / 0.17 & 1.18 & 1.17 & 1.18 & 0.80 & 0.48 \\
OG Rak 06 ZZ & sand & 0.57 / 0.29 / 0.14 & 0.82 & 0.82 & 0.82 & 0.57 & 0.33 \\
OG Rak 08 OZ & mixed & 0.56 / 0.11 / 0.33 & 1.38 & 1.38 & 1.38 & 1.30 & 1.28 \\
OG Rak 09 ZOZ & clay/peat & 0.50 / 0.17 / 0.33 & 1.10 & 1.10 & 1.10 & 1.14 & 1.21 \\
OG Rak 09 WZ & clay/peat & 0.56 / 0.11 / 0.33 & 1.05 & 1.05 & 1.05 & 1.08 & 1.14 \\
OG Rak 09 WZ 01 & clay/peat & 0.50 / 0.17 / 0.33 & 1.13 & 1.13 & 1.14 & 1.19 & 1.26 \\
KNG 002 & clay/peat & 0.67 / 0.00 / 0.33 & 1.37 & 1.37 & 1.38 & 1.39 & 1.43 \\
KNG 003 & mixed & 0.75 / 0.00 / 0.25 & 1.52 & 1.52 & 1.53 & 1.48 & 1.37 \\
KNG 004 & mixed & 0.75 / 0.00 / 0.25 & 1.39 & 1.39 & 1.40 & 1.37 & 1.23 \\
KNG 005 & clay/peat & 0.50 / 0.00 / 0.50 & 0.59 & 0.59 & 0.59 & 0.67 & 0.65 \\
KNG 006 & clay/peat & 0.40 / 0.20 / 0.40 & 0.59 & 0.60 & 0.60 & 0.67 & 0.73 \\
KNG 007 & clay/peat & 0.40 / 0.20 / 0.40 & 0.59 & 0.59 & 0.59 & 0.67 & 0.73 \\
KNG 008 & clay/peat & 0.40 / 0.20 / 0.40 & 0.56 & 0.56 & 0.56 & 0.60 & 0.85 \\
KNG 009 & sand & 0.67 / 0.33 / 0.00 & 1.44 & 1.44 & 1.45 & 1.12 & 0.81 \\
\bottomrule
\end{tabular}
\end{table}

\paragraph{Correlations.} Both correlations change $\beta$ by less than the
Monte Carlo error of these runs ($\pm 0.01$) for most walls, and by at most
0.03. A negative $\rho(K,\gamma)$ narrows the spread of $K\gamma$, which drives
the soil load. A positive $\rho(\text{MOR},E)$ couples strength to stiffness:
a stiffer pile attracts a larger moment ($M \propto E^{1/4}$), but its higher
MOR more than covers this.

\paragraph{Groundwater.} Higher groundwater adds water pressure behind the
wall, so $\beta$ falls for each type. For the best estimate the effect is
partly offset: under the higher water fewer type-1 draws stand today, so the
update moves $P(T{=}2)$ up. Walls in sand, where type 1 is likely, lose most.

\paragraph{Soil class.} Most walls have borings of all three classes within
150\,m. Walls classed as sand lose up to 0.7 in $\beta$, because half of their
nearby borings find clay or peat. Walls classed as clay or peat gain up to
0.6 (Rak 03 OZ, where half of the nearby borings find sand).

\paragraph{The two limit states.} For type 2, overturning and pile failure
occur together 4.4 times as often (median) as if they were independent: both
grow with the soil load. Because overturning is rare for type 2, the series
$P_f$ differs by less than 8\,\% from the value for independent limit states.

\subsection{One-at-a-time sensitivity}\label{app:prob-tornado}

The script \texttt{SensitivityTornado.py} (in
\texttt{scripts/StressCaclulation/}) supports Section~\ref{sec:dimension-reduction}.
For each of the 17 walls it reruns the model once as it is, twice for each of
the 18 random inputs, fixed at its 5th and at its 95th percentile, once for
each soil class, and once with six minor inputs fixed together at their
medians, with 100{,}000 draws per wall, type and run. A fixed input still uses
up its random numbers, so all other inputs get the same draws in every run and
the differences are not Monte Carlo noise. It writes
\texttt{walmuur\_tornado.csv} in \texttt{data/processed/}.
Table~\ref{tab:tornado-full} gives the change in $\beta$ (updated, 2026, for
the shell growth 2041) from the 5th to the 95th percentile. For Kromme Nieuwe
Gracht 005 to 007 the werf level is an assumed range, not a measurement error,
so it is shown separately and is not fixed in the last row. Between the wall
types, $\beta$ differs by 2.06 (median; 1.36 to 2.41), and between the
highest and lowest soil class by 0.49 (0.37 to 1.19, best estimate).

\begin{table}[htbp]
\centering
\footnotesize
\caption{Change in $\beta$ when one input moves from its 5th to its 95th
percentile: median [range] over the walls. Best estimate: both types weighted
by $P(T \mid E)$.}
\label{tab:tornado-full}
\setlength{\tabcolsep}{3pt}
\begin{tabular}{@{}l l r r r@{}}
\toprule
Input & 5th to 95th percentile & Type 1 & Type 2 & Best estimate \\
\midrule
Earth pressure coefficient $K$ & Per class\textsuperscript{a} & \textminus2.38 [\textminus3.18, \textminus1.71] & \textminus0.97 [\textminus1.23, \textminus0.68] & \textminus0.38 [\textminus2.07, +0.57] \\
Distance $a$ of the surcharge & 0.15--2.85\,m & +1.45 [+1.08, +1.57] & +1.01 [+0.73, +1.12] & +0.74 [+0.69, +1.08] \\
Soft shell $s_0$ & 2--43\,mm & \textminus0.04 [\textminus0.17, \textminus0.02] & \textminus1.50 [\textminus2.84, \textminus1.17] & \textminus0.54 [\textminus0.95, \textminus0.18] \\
Pile diameter $D_0$ & 150--210\,mm & +0.04 [+0.01, +0.18] & +1.49 [+0.92, +2.72] & +0.40 [+0.15, +0.81] \\
Werf level, assumed (KNG 005--007) & 1.03--1.57\,m NAP & \textminus0.75 [\textminus0.77, \textminus0.73] & \textminus1.04 [\textminus1.05, \textminus1.04] & +0.04 [+0.04, +0.05] \\
Model factor $\theta_R$ & 0.77--1.26 & +0.84 [+0.67, +1.22] & +0.12 [\textminus0.02, +0.34] & \textminus0.03 [\textminus0.25, +1.10] \\
Groundwater level & +0.52 to +0.88\,m NAP & \textminus0.82 [\textminus1.46, \textminus0.63] & \textminus0.18 [\textminus0.36, \textminus0.17] & +0.05 [\textminus1.32, +0.28] \\
Foundation level & \textminus0.135 to +0.135\,m & +0.66 [+0.55, +0.91] & +0.49 [+0.37, +0.52] & +0.20 [\textminus0.04, +0.85] \\
Surcharge $q$ & 4.7--9.3\,kPa & \textminus0.60 [\textminus0.70, \textminus0.56] & \textminus0.49 [\textminus0.53, \textminus0.34] & \textminus0.36 [\textminus0.55, \textminus0.31] \\
Pile spacing $s$ & 0.82--1.18\,m & \textminus0.01 [\textminus0.04, 0.00] & \textminus0.41 [\textminus0.46, \textminus0.28] & \textminus0.12 [\textminus0.25, \textminus0.04] \\
Bending strength MOR & 18.4--25.6\,MPa & +0.01 [0.00, +0.04] & +0.37 [+0.26, +0.41] & +0.12 [+0.04, +0.23] \\
Werf level, measured (14 walls) & AHN $\pm$0.08\,m & \textminus0.26 [\textminus0.31, \textminus0.22] & \textminus0.32 [\textminus0.39, \textminus0.25] & \textminus0.16 [\textminus0.28, 0.00] \\
Model factor $\theta_M$ & 0.84--1.17 & \textminus0.01 [\textminus0.03, 0.00] & \textminus0.32 [\textminus0.34, \textminus0.22] & \textminus0.10 [\textminus0.20, \textminus0.03] \\
Unit weight of masonry & 17.4--20.6\,kN/m$^3$ & +0.20 [+0.15, +0.27] & 0.00 [0.00, 0.00] & \textminus0.01 [\textminus0.10, +0.25] \\
Subgrade modulus $k$ & 4{,}585--9{,}715\,kN/m$^3$ & 0.00 [0.00, +0.02] & +0.18 [+0.13, +0.20] & +0.06 [+0.02, +0.12] \\
Modulus of elasticity $E$ & 6.0--12.0\,GPa & 0.00 [\textminus0.01, 0.00] & \textminus0.16 [\textminus0.19, \textminus0.11] & \textminus0.05 [\textminus0.10, \textminus0.02] \\
Unit weight of soil $\gamma$ & Per class & +0.13 [+0.12, +0.20] & \textminus0.11 [\textminus0.14, \textminus0.03] & +0.05 [\textminus0.05, +0.18] \\
Unit weight tamped concrete & 21.1--24.9\,kN/m$^3$ & 0.00 [0.00, 0.00] & +0.02 [0.00, +0.03] & 0.00 [0.00, +0.02] \\
Shell growth $r$ (2041) & 0.11--0.24\,mm/yr & \textminus0.01 [\textminus0.03, \textminus0.01] & \textminus0.14 [\textminus0.18, \textminus0.12] & \textminus0.06 [\textminus0.11, \textminus0.02] \\
Six minor inputs at their medians\textsuperscript{b} & -- & 0.00 [\textminus0.01, +0.01] & +0.01 [0.00, +0.03] & 0.00 [\textminus0.01, +0.01] \\
\bottomrule
\multicolumn{5}{@{}p{0.97\textwidth}@{}}{\textsuperscript{a} Sand 0.25--0.48,
mixed 0.29--0.65, clay or peat 0.42--0.81. \textsuperscript{b} Measured werf
level, unit weights of soil and tamped concrete, $E$, $k$ and shell growth.}
\end{tabular}
\end{table}

\subsection{Checks against data}\label{app:prob-data}

Table~\ref{tab:prob-data} compares the input distributions that can be
checked with the data available for the 17 walls.

\begin{table}[htbp]
\centering
\small
\caption{Input distributions against the available data, for the 17 walls.}
\label{tab:prob-data}
\begin{tabularx}{\textwidth}{@{}l >{\raggedright\arraybackslash}p{3.4cm}
  >{\raggedright\arraybackslash}X >{\raggedright\arraybackslash}p{3.6cm}@{}}
\toprule
Input & Model & Data & Verdict \\
\midrule
Werf level & $N(\text{median}, 0.05\,\text{m})$ &
  AHN points below +1.6\,m NAP: SD 0.07\,m (0.01--0.33) per wall; standard error of the mean about 0.04\,m &
  Fits the wall average once points on stairs or the street are left out \\
Groundwater & $U(+0.5, +0.9)$\,m NAP &
  BRO, $1/d^2$ over 6 wells, nearest 146--309\,m away: GLG +0.90 to +1.64, GHG +1.43 to +2.15\,m NAP &
  Data not representative of the werf; physical bound used \\
Soil class & Class of nearest boring &
  25--67\,\% (median 40\,\%) of borings within 150\,m agree &
  Weak; class should be uncertain \\
Soft shell & $U(0, 45)$\,mm, mean 22.5 &
  Amsterdam pine piles: mean 21\,mm, max.\ 45\,mm \cite{hemel2023} &
  Fits the measured range; no Utrecht data \\
Surcharge & $N(7\,\text{kPa}, 0.20)$ &
  No load data; 2\,t axle limit and a pedestrian werf &
  Not testable; on the safe side \\
\bottomrule
\end{tabularx}
\end{table}

\subsection{Monte Carlo accuracy}\label{app:prob-mc}

With $n$ draws that stand today, the coefficient of variation of the updated
estimate is
\begin{equation}
\text{CoV}(\hat P_f) = \sqrt{\frac{1 - \hat P_f}{n \hat P_f}}.
\end{equation}
The worst case is Kromme Nieuwe Gracht 009, type 2, with $\hat P_f = 0.007$
and $n = 398{,}880$: a CoV of 1.9\,\%. The standard error of $\beta$ is at most
0.01 for every wall and type.

\section{Source code}
\label{app:pm-code}

